\documentclass[10pt,a4paper]{amsart}
\usepackage[margin=19mm]{geometry}
\usepackage{amsmath,amssymb,mathtools}
\usepackage[scr=rsfs,cal=euler]{mathalfa}
\usepackage{xfrac}
\usepackage[unicode,hidelinks,bookmarks=false]{hyperref}
\newcommand{\ie}{i.e.\ }
\newcommand{\cf}{cf.\ }
\newcommand{\resp}{respectively\ }
\newcommand{\Subsection}{\subsection}
\providecommand{\nobreakdash}{\nobreak\hbox{-}}
\setlength{\emergencystretch}{2em}
\multlinegap0pt
\usepackage{bm}
\usepackage{enumerate}
\usepackage{amssymb}
\newtheorem{theorem}{Theorem}
\title{KKT conditions for optimization with generalized invex fuzzy functions}
\author{Ville Rinne, Yury Nikulin, Marko M. M\"akel\"a}
\date{Source: arXiv:2603.00125v1 (CC BY 4.0)}
\begin{document}
\maketitle

\noindent\emph{Source and licence.} KKT conditions for optimization with generalized invex fuzzy functions. Version arXiv:2603.00125v1 (CC BY 4.0), distributed by arXiv under the Creative Commons Attribution 4.0 International licence, http://creativecommons.org/licenses/by/4.0/, as recorded in the arXiv licence metadata for this exact version and shown on the abstract page https://arxiv.org/abs/2603.00125v1. Full licence notice: \texttt{licenses/arxiv-2603.00125-CC-BY-4.0.txt}.\par\smallskip

\noindent\textit{Editorial scope.} Counted unit: the article's theorem 4.3.1 (source0.tex lines 1178--1188). The article's complete original proof of it is delivered with it (source0.tex lines 1189--1208, one \texttt{proof} environment of 20 lines). No mathematics is written, repaired or re-derived: the statement and the proof are the article's own lines, in the article's own notation, and no retained line is substituted or re-flowed.  No further source-local context is needed: every cross-reference of every retained line resolves inside the two delivered spans.  Nothing was omitted from any retained span, so the package carries no omission record.  No retained line contains a citation, so the wrapper carries no bibliography and no bibliography entry.  No retained line contains a figure environment, an image-inclusion command or drawing code, so no image is delivered and the package declares no asset dependency.  Editorial formatting only, in addition: an AMS \texttt{amsart} preamble that restates the article's own declarations and macros used by the retained lines, namely the environments \texttt{theorem} on separate counters, together with the standard packages those lines need.  The retained environments print in this document as Theorem 1 in order of appearance, whereas the article prints the same items with the numbering of its own full document.  The complete native source of the article as distributed in the arXiv e-print for this exact version is bundled unchanged as \texttt{sources/195470/source0.tex}.
\begin{theorem}\label{4.3.1} Let $\boldsymbol u$ be a feasible solution of the considered nonsmooth FOP (all functions $f$, $g$, and $h$ are LLC).
Let $J(\boldsymbol{u}) = \{j\!:\,g_j(\boldsymbol{u}) = 0\}$. For some $\alpha\in[0,1]$, assume that there exist $\lambda_1( \alpha) > 0,\  \lambda_2( \alpha) > 0,\  {\boldsymbol \mu}( \alpha) \in \mathbb  R^{m}$ with
${\boldsymbol\mu}(\alpha)\ge \boldsymbol 0$, and $ {\boldsymbol \vartheta}( \alpha) \in \mathbb R^r$ such that KKT optimality conditions
$$\boldsymbol 0 \in  \lambda_1( \alpha)\partial_Cf^L(\boldsymbol u,  \alpha) +  \lambda_2( \alpha)\partial_Cf^R(\boldsymbol u,  \alpha)$$
$$+ \sum_{j\in J(\boldsymbol{u})}  \mu_j( \alpha)\partial_Cg_j(\boldsymbol u) + \sum_{k=1}^r  \vartheta_k( \alpha)\partial_Ch_k(\boldsymbol u),$$
$$ \mu_j( \alpha)g_j(\boldsymbol u) = 0,\ j\in J$$
hold. If $$( \lambda_1( \alpha) f^L(\boldsymbol x,  \alpha),  \lambda_2( \alpha) f^R(\boldsymbol x,  \alpha))$$ is V-pseudoinvex at $\boldsymbol u$ on $X$ with respect to $\boldsymbol \eta$, 
%$$( \mu_1( \alpha) g_1(\boldsymbol x),\dots, \mu_m( \alpha) g_m(\boldsymbol x))$$ 
$$( \mu_j( \alpha) g_j(\boldsymbol x))_{j \in J(\boldsymbol{u})}$$
and $$( \vartheta_1( \alpha) h_1(\boldsymbol x),\dots,  \vartheta_r( \alpha) h_r(\boldsymbol x))$$ are V-quasiinvex at $\boldsymbol u$ on $X$ with respect to the same $\boldsymbol \eta$, then $\boldsymbol u$ is a weakly $\tilde f$-nondominated solution of the FOP.
\end{theorem}

\begin{proof} By assumption, $\boldsymbol u$ is a feasible solution of the nonsmooth FOP such that for some $ \alpha \in [0,1]$ there exist $\lambda_1( \alpha) > 0$, $\lambda_2( \alpha) > 0$, $\boldsymbol { \mu}( \alpha) \in \mathbb R^{m},\ \boldsymbol{ \mu}( \alpha) \geq \boldsymbol 0$, $\boldsymbol { \vartheta}( \alpha) \in \mathbb R^r$ such that the KKT optimality conditions hold. By assumption, all functions involved in the FOP are LLC. Therefore, all functions constituting the associated scalarized problem $\boldsymbol {P_{\alpha}({\boldsymbol \lambda)}}$ are also LLC. By the KKT conditions, we know that there exist $\boldsymbol \xi_L \in \partial_C f^L(\boldsymbol u,  \alpha),\ \boldsymbol \xi_R \in \partial_C f^R(\boldsymbol u,  \alpha)$, $\{\boldsymbol \varsigma_j | \boldsymbol \varsigma_j \in \partial_C g_j(\boldsymbol u),\ j \in J(\boldsymbol{u})\}$ and $\{\boldsymbol \zeta_k | \boldsymbol \zeta_k \in \partial_C h_k(\boldsymbol u),\ k=1,\dots,r\}$ such that
$$ \lambda_1( \alpha) \boldsymbol \xi_L +  \lambda_2( \alpha) \boldsymbol \xi_R + \sum_{j \in J(\boldsymbol{u})} \mu_j( \alpha) \boldsymbol \varsigma_j + \sum_{k=1}^r \vartheta_k( \alpha) \boldsymbol \zeta_k = \boldsymbol 0.$$
Therefore, for all feasible $\boldsymbol x\in X$, we have 
$$ \lambda_1( \alpha) \boldsymbol \xi_{L}^T \boldsymbol \eta(\boldsymbol x; \boldsymbol u) +  \lambda_2( \alpha) \boldsymbol \xi_{R}^T \boldsymbol \eta(\boldsymbol x; \boldsymbol u) +$$ $$\sum_{j \in J(\boldsymbol{u})} \mu_j( \alpha) \boldsymbol \varsigma_j^T \boldsymbol \eta(\boldsymbol x; \boldsymbol u)+\sum_{k=1}^r \vartheta_k( \alpha) \boldsymbol \zeta^T_k \boldsymbol \eta(\boldsymbol x; \boldsymbol u) = 0.$$
Since $ \mu_j( \alpha) g_j(\boldsymbol u) = 0$ and $ \mu_j( \alpha) g_j(\boldsymbol x) \leq 0$ for any feasible $\boldsymbol x$ and $j\in J$, we have $$ \mu_j( \alpha) g_j(\boldsymbol x) \leq  \mu_j( \alpha) g_j(\boldsymbol u).$$
Due to the definition of $V$-quasiinvexity, there exist $\beta_j : X \times X \rightarrow \mathbb R_{> \boldsymbol 0},\ j=1,\dots,m$. Since $\beta_j(\boldsymbol x; \boldsymbol u) > 0$, we have 
$$\sum_{j \in J(\boldsymbol{u})}  \mu_j( \alpha) \beta_j(\boldsymbol x; \boldsymbol u) g_j(\boldsymbol x) \leq \sum_{j \in J(\boldsymbol{u})}  \mu_j( \alpha) \beta_j(\boldsymbol x; \boldsymbol u) g_j(\boldsymbol u).$$
Then, by the V-quasiinvexity condition, we get
$$\sum_{j \in J(\boldsymbol{u})}  \mu_j( \alpha) \boldsymbol \varsigma_j^T \boldsymbol \eta(\boldsymbol x; \boldsymbol u) \leq 0,\ \forall\ \boldsymbol \varsigma_j \in \partial_C g_j(\boldsymbol u).$$
Since $ \vartheta_k( \alpha) h_k(\boldsymbol x) =  \vartheta_k( \alpha) h_k(\boldsymbol u) = 0$ for any feasible $\boldsymbol x$, we can show that similar nonpositivity constraints are true for the equality constraints functions $h_k$. Thus, we have
$$ \lambda_1( \alpha) \boldsymbol \xi_{L}^T \boldsymbol \eta(\boldsymbol x; \boldsymbol u) +  \lambda_2( \alpha) \boldsymbol \xi_{R}^T \boldsymbol \eta(\boldsymbol x; \boldsymbol u) \geq 0.$$
Then, again, by the V-pseudoinvexity condition, there exists a vector $\boldsymbol \beta\in \mathbb R^2_{> \boldsymbol 0}$ with components $\beta_i: X \times X \rightarrow \mathbb R_{> \boldsymbol 0}$, $i=1,2$, such that for all feasible $\boldsymbol x\in X$,  we have

$$\beta_1(\boldsymbol x, \boldsymbol u)  \lambda_1( \alpha) f^L(\boldsymbol x,  \alpha) + \beta_2(\boldsymbol x, \boldsymbol u)  \lambda_2( \alpha) f^R(\boldsymbol x,  \alpha)$$
$$\geq \beta_1(\boldsymbol x, \boldsymbol u)  \lambda_1( \alpha) f^L(\boldsymbol u,  \alpha) + \beta_2(\boldsymbol x, \boldsymbol u)  \lambda_2( \alpha) f^R(\boldsymbol u,  \alpha).$$
We assume that $\boldsymbol u$ is not a Pareto optimal point for $(f^L,f^R)$. Then there exists $\boldsymbol x^*\in X$ such that $$f^L(\boldsymbol x^*) \leq f^L(\boldsymbol u)\text{ and } f^R(\boldsymbol x^*) \leq f^R(\boldsymbol u),$$ as well as $$f^L(\boldsymbol x^*) < f^L(\boldsymbol u)\text{ or } f^R(\boldsymbol x^*) < f^R(\boldsymbol u).$$ Then we have
$$\beta_1(\boldsymbol x, \boldsymbol u)  \lambda_1( \alpha) f^L(\boldsymbol x^*,  \alpha) + \beta_2(\boldsymbol x, \boldsymbol u)  \lambda_2( \alpha) f^R(\boldsymbol x^*,  \alpha)$$
$$< \beta_1(\boldsymbol x, \boldsymbol u)  \lambda_1( \alpha) f^L(\boldsymbol u,  \alpha) + \beta_2(\boldsymbol x, \boldsymbol u)  \lambda_2( \alpha) f^R(\boldsymbol u,  \alpha).$$
This is a contradiction. Hence, $\boldsymbol u$ is a Pareto optimal solution for VMP$_{\alpha}$. Since $\alpha$ is fixed, $\boldsymbol u$ is a weakly $\tilde f$-nondominated solution of the FOP. 
\end{proof}

\end{document}
